% Appendix C: Practitioner Checklist Cards
% Condensed, scannable reformatting of Section 07 (Guidance Artifact)
% Each card fits approximately one page

\section[Appendix - Practitioner Checklist Cards]{Practitioner Checklist Cards}\label{app:checklists}

This appendix condenses the guidance artifact (Section~\ref{sec:artifact}) into
twelve scannable checklist cards. Each card traces directly to a specific
subsection of the artifact; no new content is introduced. Practitioners may
photocopy or print individual cards for use during Reference Architecture (RA)
development sessions. For Cards~1--7, practice descriptions and verification
patterns carry Delphi consensus; experience indicators (Novice/Advanced) are
researcher proposals.

%% ================================================================
%% Card 1: System/Codebase Analysis
%% ================================================================
\clearpage
\subsection*{Card 1: System/Codebase Analysis \hfill \textit{[Directive]}}

\begin{tabular}{|p{7cm}|p{6cm}|}
    \hline
    \textbf{Consensus:} $M = 4.00$, $SD = 0.82$ & \textbf{Provenance:} Delphi consensus \\
    \hline
\end{tabular}

\vspace{0.5em}
\textbf{When to Use}
\begin{itemize}
    \item Phase~2 (Current State Analysis) of the RA lifecycle
    \item Legacy system comprehension, dependency mapping, cross-repository analysis
\end{itemize}

\textbf{Key Steps}
\begin{enumerate}
    \item Feed existing system artifacts (code, configs, diagrams) into GenAI for analysis
    \item Delegate information gathering to GenAI; retain architectural judgment
    \item Cross-validate GenAI-generated descriptions against source systems directly
    \item Verify that referenced components, services, and dependencies exist in the actual system
\end{enumerate}

\textbf{Verification Checklist}
\begin{itemize}[label=\checkbox]
    \item Cross-validate system descriptions against source systems (Trust level: Highest)
    \item Confirm referenced folders, components, and services exist physically
    \item Verify dependency graphs against actual system topology
    \item Spot-check representative samples of information compression outputs
\end{itemize}

\begin{tabular}{|p{6.5cm}|p{6.5cm}|}
    \hline
    \textbf{Novice} & \textbf{Advanced} \\
    \hline
    Start with well-documented codebases where verification is straightforward & Extend to legacy systems and cross-repository analysis where GenAI's compression value is highest \\
    \hline
\end{tabular}

\vspace{0.5em}
\textit{External corroboration: Section~\ref{subsec:disc-triangulation}}

%% ================================================================
%% Card 2: Research & Exploration
%% ================================================================
\clearpage
\subsection*{Card 2: Research \& Exploration \hfill \textit{[Contextual]}}

\begin{tabular}{|p{7cm}|p{6cm}|}
    \hline
    \textbf{Consensus:} $M = 3.60$, $SD = 1.17$ & \textbf{Provenance:} Delphi consensus \\
    \hline
\end{tabular}

\vspace{0.5em}
\textbf{When to Use}
\begin{itemize}
    \item Phase~4 (Target State Design) of the RA lifecycle
    \item Exploring design alternatives, technology options, and architectural trade-offs
\end{itemize}

\textbf{Key Steps}
\begin{enumerate}
    \item Use GenAI as a research assistant for information gathering
    \item Scope research questions clearly before prompting
    \item Treat GenAI research outputs as hypotheses requiring confirmation
    \item Verify claims, statistics, and technology recommendations against primary sources
\end{enumerate}

\textbf{Verification Checklist}
\begin{itemize}[label=\checkbox]
    \item Cross-validate specific claims against primary sources (heuristic $M = 4.57$; Trust level: Medium-Low)
    \item Verify 60--70\% of specific claims against primary sources (P2, R1 Workbook)
    \item Confirm technology recommendations against official documentation
    \item Check statistics and benchmarks for currency and accuracy
\end{itemize}

\begin{tabular}{|p{6.5cm}|p{6.5cm}|}
    \hline
    \textbf{Novice} & \textbf{Advanced} \\
    \hline
    Focus on well-scoped research questions where outputs can be verified against established sources & Leverage multi-model workflows and iterative refinement cycles \\
    \hline
\end{tabular}

\vspace{0.5em}
\textit{External corroboration: Section~\ref{subsec:disc-triangulation}}

%% ================================================================
%% Card 3: Adversarial Review / Stress Testing
%% ================================================================
\clearpage
\subsection*{Card 3: Adversarial Review / Stress Testing \hfill \textit{[Contextual]}}

\begin{tabular}{|p{7cm}|p{6cm}|}
    \hline
    \textbf{Consensus:} $M = 3.60$, $SD = 0.97$ & \textbf{Provenance:} Delphi consensus \\
    \hline
\end{tabular}

\vspace{0.3em}
\noindent\textit{Note: Contextual despite low SD ($< 1.0$) because practice value depends on practitioner domain expertise.}

\vspace{0.5em}
\textbf{When to Use}
\begin{itemize}
    \item When challenging, critiquing, or stress-testing architectural decisions
    \item Requires domain expertise as a prerequisite
\end{itemize}

\textbf{Key Steps}
\begin{enumerate}
    \item Select a stress-testing technique (see list below)
    \item Actively elicit critique; do not accept GenAI's default collaborative stance
    \item Evaluate GenAI-generated critiques against your own domain knowledge
    \item Recognize that GenAI cannot model strategic adversarial behaviour across organizational boundaries
\end{enumerate}

\textbf{Stress-Testing Techniques}
\begin{itemize}
    \item AI Sparring Partner: iterative adversarial dialogue challenging design rationale
    \item Devil's Advocate Prompting: instruct GenAI to argue against the proposed architecture
    \item Decision Stress Testing: identify blind spots, steelman alternatives, conduct pre-mortem
    \item AI Pre-Mortem: ``The architecture already failed; why?''
    \item Skeptical Reviewer: persona critique simulating a senior architect review
    \item Architectural Gap Analysis: completeness analysis against standards
    \item Epistemic Stress Testing: challenge confidence levels, not decisions themselves
\end{itemize}

\textbf{Verification Checklist}
\begin{itemize}[label=\checkbox]
    \item Confirm domain expertise threshold before relying on adversarial outputs ($M = 3.86$)
    \item Distinguish genuine design weaknesses from pattern-matching artifacts
    \item Validate contradiction seeding results against known architectural constraints ($M = 3.33$)
\end{itemize}

\begin{tabular}{|p{6.5cm}|p{6.5cm}|}
    \hline
    \textbf{Novice} & \textbf{Advanced} \\
    \hline
    Start with structured techniques (AI Pre-Mortem, Skeptical Reviewer) that provide clear prompting templates & Develop intuition for when GenAI critique reflects genuine weaknesses versus pattern-matching artifacts \\
    \hline
\end{tabular}

\vspace{0.5em}
\textit{External corroboration: Section~\ref{subsec:disc-triangulation}}

%% ================================================================
%% Card 4: Governance/Compliance Checking
%% ================================================================
\clearpage
\subsection*{Card 4: Governance/Compliance Checking \hfill \textit{[Adaptive]}}

\begin{tabular}{|p{7cm}|p{6cm}|}
    \hline
    \textbf{Consensus:} $M = 3.50$, $SD = 1.27$ & \textbf{Provenance:} Delphi consensus \\
    \hline
\end{tabular}

\vspace{0.5em}
\textbf{When to Use}
\begin{itemize}
    \item Verifying architectural artifacts against governance standards and compliance requirements
    \item Checking organizational policies; adapt to your specific governance structures
\end{itemize}

\textbf{Key Steps}
\begin{enumerate}
    \item Scan documents against regulatory requirements using GenAI
    \item Verify completeness and accuracy manually in all compliance contexts
    \item Apply line-by-line verification for compliance-sensitive outputs ($M = 4.20$)
    \item Recognize that compliance errors carry disproportionate organizational risk
\end{enumerate}

\textbf{Verification Checklist}
\begin{itemize}[label=\checkbox]
    \item Verify every compliance-related detail line-by-line (Trust level: Very Low)
    \item Confirm GenAI output against the actual regulatory text or standard
    \item Check for completeness gaps that GenAI may have omitted
    \item Validate that organizational-specific policies are correctly represented
    \item Verify output structure and formatting against expected standards ($M = 2.83$)
\end{itemize}

\begin{tabular}{|p{6.5cm}|p{6.5cm}|}
    \hline
    \textbf{Novice} & \textbf{Advanced} \\
    \hline
    Use GenAI for compliance checking only in low-risk contexts initially & Integrate governance checks into automated pipelines where organizational maturity permits \\
    \hline
\end{tabular}

\vspace{0.5em}
\textit{External corroboration: Section~\ref{subsec:disc-triangulation}}

%% ================================================================
%% Card 5: First Draft & Documentation
%% ================================================================
\clearpage
\subsection*{Card 5: First Draft \& Documentation \hfill \textit{[Contextual]}}

\begin{tabular}{|p{7cm}|p{6cm}|}
    \hline
    \textbf{Consensus:} $M = 3.10$, $SD = 1.10$ & \textbf{Provenance:} Delphi consensus \\
    \hline
\end{tabular}

\vspace{0.5em}
\textbf{When to Use}
\begin{itemize}
    \item Producing initial drafts of architectural documentation (solution designs, ADRs, technical specifications, assessment reports)
    \item Apply the Compression Not Generation framing: GenAI compresses existing information into structured documents
\end{itemize}

\textbf{Key Steps}
\begin{enumerate}
    \item Provide existing information (meeting notes, system data, design inputs) as context
    \item Use GenAI as a First Draft Generator; treat outputs as starting points, not finished artifacts
    \item Cross-check all references, links, and factual claims in generated documentation
    \item Apply architectural judgment to refine structure and content
\end{enumerate}

\textbf{Verification Checklist}
\begin{itemize}[label=\checkbox]
    \item Cross-check all references and links for accuracy (Trust level: High)
    \item Verify factual claims and benchmarks against primary sources
    \item Check for hallucinated references or non-existent sources
    \item Confirm document structure aligns with organizational templates and standards
\end{itemize}

\begin{tabular}{|p{6.5cm}|p{6.5cm}|}
    \hline
    \textbf{Novice} & \textbf{Advanced} \\
    \hline
    Start with well-structured document types (meeting minutes, status reports) before progressing to ADRs & Develop metaprompt libraries and multi-stage generation workflows \\
    \hline
\end{tabular}

\vspace{0.5em}
\textit{External corroboration: Section~\ref{subsec:disc-triangulation}}

%% ================================================================
%% Card 6: Stakeholder Communication
%% ================================================================
\clearpage
\subsection*{Card 6: Stakeholder Communication \hfill \textit{[Adaptive]}}

\begin{tabular}{|p{7cm}|p{6cm}|}
    \hline
    \textbf{Consensus:} $M = 3.10$, $SD = 1.37$ & \textbf{Provenance:} Delphi consensus \\
    \hline
\end{tabular}

\vspace{0.5em}
\textbf{When to Use}
\begin{itemize}
    \item Translating architectural content for different audiences (executive summaries, team-specific views, non-technical explanations)
    \item Synthesizing stakeholder inputs into structured architectural artifacts
\end{itemize}

\textbf{Key Steps}
\begin{enumerate}
    \item Define the target audience and communication purpose before prompting
    \item Inject organizational context (team dynamics, political constraints, historical decisions) manually
    \item Validate that GenAI outputs reflect organizational reality, not just technical logic
    \item Review for audience appropriateness before distribution
\end{enumerate}

\textbf{Verification Checklist}
\begin{itemize}[label=\checkbox]
    \item Validate organizational context manually before communicating outputs (Trust level: Lowest for organizational judgment)
    \item Check that outputs are not ``technically logical but organisationally nonsensical'' (P3, R1 Workbook)
    \item Verify that audience-specific framing preserves architectural intent
    \item Confirm that political and interpersonal nuances are appropriately handled
    \item Simulate likely stakeholder reactions before presenting GenAI-assisted content ($M = 3.00$)
\end{itemize}

\begin{tabular}{|p{6.5cm}|p{6.5cm}|}
    \hline
    \textbf{Novice} & \textbf{Advanced} \\
    \hline
    Limit GenAI use to summarization tasks (meeting minutes, document distillation) where source material constrains outputs & Develop audience-specific prompting strategies and RAG-based approaches over organizational documentation \\
    \hline
\end{tabular}

\vspace{0.5em}
\textit{External corroboration: Section~\ref{subsec:disc-triangulation}}

%% ================================================================
%% Card 7: Deep Workflow Integration
%% ================================================================
\clearpage
\subsection*{Card 7: Deep Workflow Integration \hfill \textit{[Adaptive]}}

\begin{tabular}{|p{7cm}|p{6cm}|}
    \hline
    \textbf{Consensus:} $M = 2.89$, $SD = 1.17$ (one abstention) & \textbf{Provenance:} Delphi consensus \\
    \hline
\end{tabular}

\vspace{0.5em}
\textbf{When to Use}
\begin{itemize}
    \item Embedding GenAI into daily workflow through IDE integration, agent context files, custom hooks, MCP servers, and multi-agent orchestration
    \item Low rating reflects adoption barriers (cognitive load, organizational resistance), not low practice value
\end{itemize}

\textbf{Key Steps}
\begin{enumerate}
    \item Begin with single-tool integration (IDE assistant) as a foundation
    \item Build trust incrementally; expand GenAI autonomy from low-stakes to higher-stakes tasks ($M = 4.29$)
    \item Progress to custom context files and agent orchestration as calibration develops
    \item Reduce human oversight touchpoints only after sustained positive calibration
\end{enumerate}

\textbf{Verification Checklist}
\begin{itemize}[label=\checkbox]
    \item Apply incremental trust building before expanding automation scope (Trust level: varies by task)
    \item Verify automated workflow outputs at regular intervals
    \item Confirm that reduced oversight touchpoints are justified by calibration evidence
    \item Maintain rollback capability for agentic workflow outputs
\end{itemize}

\begin{tabular}{|p{6.5cm}|p{6.5cm}|}
    \hline
    \textbf{Novice} & \textbf{Advanced} \\
    \hline
    Begin with single-tool integration (IDE assistant) before progressing to custom context files and agent orchestration & Design multi-agent architectures with specialized agents for different architectural tasks \\
    \hline
\end{tabular}

\vspace{0.5em}
\textit{External corroboration: Section~\ref{subsec:disc-triangulation}}

%% ================================================================
%% Card 8: Trust Calibration
%% ================================================================
\clearpage
\subsection*{Card 8: Trust Calibration \hfill \textit{Cross-Cutting}}

\begin{tabular}{|p{7cm}|p{6cm}|}
    \hline
    \textbf{Cross-validation heuristic:} $M = 4.57$, $SD = 0.53$ & \textbf{Provenance:} Researcher synthesis (panel validated heuristics) \\
    \hline
\end{tabular}

\vspace{0.5em}
\textbf{Trust Spectrum} (from Table~\ref{tab:trust-spectrum})

\begin{tabular}{|p{2.5cm}|p{5cm}|p{5cm}|}
    \hline
    \textbf{Trust Level} & \textbf{Task Type} & \textbf{Verification Approach} \\
    \hline
    Highest & Information compression: summaries, boilerplate, dependency graphs & Spot-check representative samples \\
    \hline
    High & Documentation drafts, concept explanation, test generation & Review structure and key claims \\
    \hline
    Medium & Code refactoring, common patterns, meeting synthesis & Verify against organizational standards \\
    \hline
    Medium-Low & Design trade-off exploration, technology recommendations & Cross-validate with primary sources \\
    \hline
    Low & Architecture decisions, novel business logic & Independent domain expert review \\
    \hline
    Very Low & Security-critical code, compliance assessment & Line-by-line verification \\
    \hline
    Lowest & Organizational judgment: politics, team dynamics, legal & Retain entirely; do not delegate \\
    \hline
\end{tabular}

\vspace{0.5em}
\textbf{Calibration Checklist}
\begin{itemize}[label=\checkbox]
    \item Match each GenAI task to its trust level before accepting output
    \item Apply the ``junior colleague'' mental model: supervisory attention for all outputs
    \item Recalibrate trust upward when evidence warrants (bidirectional calibration)
    \item Build domain expertise first; trust calibration is teachable but requires domain knowledge as a foundation
    \item Learn GenAI capability boundaries through direct experience with errors (frontier learning pattern)
\end{itemize}

\vspace{0.5em}
\textit{Source: Section~\ref{subsubsec:trust-framework}}

%% ================================================================
%% Card 9: Context Engineering
%% ================================================================
\clearpage
\subsection*{Card 9: Context Engineering for EA \hfill \textit{Cross-Cutting}}

\begin{tabular}{|p{7cm}|p{6cm}|}
    \hline
    \textbf{Context management heuristic:} $M = 4.14$, $SD = 0.90$ & \textbf{Provenance:} Researcher synthesis (panel validated context management) \\
    \hline
\end{tabular}

\vspace{0.5em}
\textbf{Agent Context Files}
\begin{itemize}
    \item Encode architectural standards, design principles, and organizational constraints in project-level files (e.g., CLAUDE.md, AGENTS.md)
    \item These files function as both human-readable documentation and machine-executable governance
\end{itemize}
\begin{itemize}[label=\checkbox]
    \item Create and maintain agent context files for each architectural project
    \item Review and update context files regularly (empirical benchmark: every one to three days)
    \item Encode architectural standards, design principles, and organizational constraints in each context file
\end{itemize}

\vspace{0.5em}
\textbf{Strategic Information Ordering}
\begin{itemize}
    \item Mitigate the ``lost in the middle'' effect: language models attend less to information in the middle of long contexts
\end{itemize}
\begin{itemize}[label=\checkbox]
    \item Place critical architectural constraints at the beginning and end of context
    \item Use modular rules files (separate documents for different architectural concerns)
\end{itemize}

\vspace{0.5em}
\textbf{Context Scope Management}
\begin{itemize}
    \item Include: architectural principles, relevant standards, system-specific technical context, decision history (ADRs)
    \item Exclude: organizational politics, exhaustive historical documentation, ambiguous or contradictory requirements
\end{itemize}
\begin{itemize}[label=\checkbox]
    \item Verify that context includes relevant constraints and excludes noise
    \item Confirm that organizational politics and interpersonal dynamics are handled by the practitioner, not delegated to GenAI
\end{itemize}

\vspace{0.5em}
\textit{Source: Section~\ref{subsubsec:context-engineering}}

%% ================================================================
%% Card 10: Anti-Atrophy Strategies
%% ================================================================
\clearpage
\subsection*{Card 10: Anti-Atrophy Strategies \hfill \textit{Cross-Cutting}}

\begin{tabular}{|p{7cm}|p{6cm}|}
    \hline
    \textbf{``Skills used less'' significance:} $M = 3.90$ & \textbf{Provenance:} Researcher proposal (panel validated atrophy concern) \\
    \hline
\end{tabular}

\vspace{0.5em}
\textbf{The Competency Decay Paradox}

The skills being automated (cold-start writing, exhaustive manual search, cover-to-cover documentation reading) are precisely the skills through which practitioners built their expertise. Three deliberate friction practices mitigate this risk.

\vspace{0.5em}
\textbf{Friction Practices}
\begin{itemize}[label=\checkbox]
    \item \textbf{AI-Free Design Sessions:} Periodically conduct architectural work without GenAI assistance to maintain foundational skills. \textit{Frequency: schedule regular sessions (analogous to manual flight training despite autopilot capabilities).}
    \item \textbf{Attempt Independently First:} Engage with a problem for 15--30 minutes before consulting GenAI. This ensures the practitioner develops their own mental model before seeing an AI-generated one. \textit{Frequency: every new design problem.}
    \item \textbf{Pair Programming Mindset:} Treat GenAI interaction as collaborative work rather than delegation. The practitioner drives architectural thinking; GenAI handles information retrieval and expression. \textit{Frequency: every GenAI session.}
\end{itemize}

\vspace{0.5em}
\textbf{Self-Assessment}
\begin{itemize}[label=\checkbox]
    \item Can you still perform core architectural tasks without GenAI assistance?
    \item Do you form your own mental model before consulting GenAI?
    \item Are you actively driving architectural thinking during GenAI sessions?
    \item Do you track GenAI failures to refine your trust calibration? (learn through failure, $M = 3.40$)
\end{itemize}

\vspace{0.5em}
\textit{Source: Section~\ref{subsubsec:anti-atrophy}}

%% ================================================================
%% Card 11: Governance Adaptation
%% ================================================================
\clearpage
\subsection*{Card 11: Governance Adaptation \hfill \textit{Organizational}}

\begin{tabular}{|p{7cm}|p{6cm}|}
    \hline
    \textbf{Evidence base:} Thinner Delphi evidence than practice catalog & \textbf{Provenance:} Mixed (see per-theme indicators) \\
    \hline
\end{tabular}

\vspace{0.5em}
\textbf{Gate-to-Guardrail Transition} \textit{(Researcher synthesis)}
\begin{itemize}[label=\checkbox]
    \item Assess whether periodic architectural review boards (ARBs) keep pace with GenAI-accelerated artifact production
    \item Consider transitioning from periodic gate reviews to continuous guardrail enforcement
    \item Encode architectural constraints for automatic checking rather than episodic evaluation
    \item Recognize that inside-out governance rated low ($M = 2.00$); most organizations have not yet confronted this shift formally
\end{itemize}

\vspace{0.5em}
\textbf{Architecture-as-Code Enforcement} \textit{(Researcher proposal)}
\begin{itemize}[label=\checkbox]
    \item Encode architectural constraints as machine-readable rules (TypeScript validators, fitness functions, CI/CD checks)
    \item Enable GenAI agents to read and apply architectural rules before generating code
    \item Implement executable ADRs where architectural decisions become automated rules in build pipelines
    \item Bridge Governance/Compliance practice with Deep Workflow Integration
\end{itemize}

\vspace{0.5em}
\textbf{AI Disclosure and Provenance} \textit{(Researcher synthesis)}
\begin{itemize}[label=\checkbox]
    \item Establish AI disclosure policies that flag AI-generated components for calibrated review
    \item Implement provenance tracking for compliance-sensitive domains
    \item Address the ``Ghost Architect'' dynamic: identify which components were AI-generated or AI-augmented
\end{itemize}

\vspace{0.5em}
\textit{Source: Section~\ref{subsec:governance-adaptation}}

%% ================================================================
%% Card 12: Maturity Model
%% ================================================================
\clearpage
\subsection*{Card 12: Maturity Model \hfill \textit{Diagnostic}}

\begin{tabular}{|p{7cm}|p{6cm}|}
    \hline
    \textbf{Progression assumption:} $M = 4.71$ (Delphi consensus) & \textbf{Provenance:} Researcher proposal (progression assumption has Delphi consensus) \\
    \hline
\end{tabular}

\vspace{0.5em}
\textbf{Stage Diagnostic} (from Table~\ref{tab:maturity-model})

\begin{tabular}{|c|p{2.8cm}|p{4.5cm}|p{4.5cm}|}
    \hline
    \textbf{Stage} & \textbf{Label} & \textbf{Characteristics} & \textbf{Observable Indicators} \\
    \hline
    1 & Experimental & Ad-hoc GenAI use for individual tasks; chat-based interaction; no structured verification & Occasional use for research or drafting; no shared practices; trust calibration through trial and error \\
    \hline
    2 & Systematic Personal & Structured personal workflows; context engineering basics; deliberate verification & Personal prompt libraries; agent context files; anti-atrophy practices; consistent trust spectrum application \\
    \hline
    3 & Team-Level Integration & Shared practices across team; team-level agent context files; Architecture-as-Code adoption & Shared context files in version control; team verification standards; collective decision capture workflows \\
    \hline
    4 & Organizational Embedding & Governance adaptation; enterprise context engineering; agentic workflows & Automated architectural guardrails; AI disclosure policies; multi-agent orchestration; enterprise MCP infrastructure \\
    \hline
\end{tabular}

\vspace{0.5em}
\textbf{Key Principles}
\begin{itemize}[label=\checkbox]
    \item Use this model as diagnostic (locating current practice), not prescriptive (mandating progression)
    \item Recognize that progression requires sustained engagement; short experience cannot substitute for long experience ($M = 2.86$)
    \item Accept that advancement may occur at different rates across practice categories (per-category measurement $M = 3.86$)
    \item Acknowledge that organizational context modulates pace but does not replace experience ($M = 3.86$)
\end{itemize}

\vspace{0.5em}
\textit{Source: Section~\ref{subsec:maturity-model}}
